\section{Background}
\label{sec:background}

\subsection{Frames, Budgets and Hitches}

A game renders a frame and hands it to the display through a present call,
\texttt{vkQueuePresentKHR} in Vulkan. The interval between consecutive
presents is the frame time, the same quantity PresentMon reports as
\texttt{MsBetweenPresents}~\cite{presentmon} and overlays such as
MangoHud~\cite{mangohud} plot. A frame has a \emph{budget}, normally the
display period or the game's frame cap: 16.7\,ms at 60\,Hz. A \emph{hitch}
is a frame that exceeds its budget by enough to be seen.

Two properties of real games complicate the picture. Under a FIFO (vsync)
present mode the display paces the game, so a healthy frame spends much of
its time blocked inside or around the present call; that wait is pacing,
not a stall. And under Proton, the thread that calls
\texttt{vkQueuePresentKHR} is usually not the game's render thread but a
translation-layer thread such as DXVK's submission thread, a few hand-offs
away from game code. \todo{Verify the DXVK and vkd3d-proton present threads
on the evaluation stack; cite source.}

\paragraph{Loader bypass} The Vulkan loader exports \texttt{vkQueuePresentKHR},
but applications that resolve device functions through
\texttt{vkGetDeviceProcAddr} receive a pointer into the first layer or the
driver and never call the export~\cite{vulkan-loader}. On our development
machine \texttt{/usr/bin/vkcube} imports no \texttt{vk*} symbol at all.
A tool that places a uprobe on the loader's symbol therefore counts zero
presents for exactly the applications that matter. An implicit layer, in
contrast, sits in the device dispatch chain however the application obtained
its pointers.

\subsection{Off-CPU Time and Wakeup Chains}

A thread that is not running is either \emph{runnable} (preempted, or woken
and waiting for a CPU) or \emph{blocked} (sleeping until an event). Linux
exposes both transitions as tracepoints: \texttt{sched\_switch} fires when a
CPU switches from one task to another and reports the outgoing task's state;
\texttt{sched\_waking} fires when a task is woken, on the waker's CPU and in
the waker's context. Off-CPU analysis~\cite{gregg-offcpu} times the blocked
intervals and attributes them to the blocking stack; off-wake and chain-graph
analysis~\cite{gregg-chaingraph} additionally records who woke the blocked
thread, and who woke that waker. Gregg's 2016 chain-graph prototype was
limited by the eBPF of the time to kernel stacks and short chains, and noted
that latency outliers were not yet addressed.

Chains are also the basis of critical-path analysis over kernel traces:
Giraldeau and Dagenais recover the wait dependencies of a task recursively
across threads and hosts~\cite{giraldeau2016wait}, wPerf builds a wait-for
graph to find bottleneck waiting events~\cite{zhou2018wperf}, and Perfetto
ships an offline critical-path library over its thread-state
tables~\cite{perfetto-critpath}. All of these operate offline, over a full
trace.

\paragraph{Interrupt context} \texttt{sched\_waking} runs in the waker's context,
but when the wakeup comes from an interrupt, \texttt{current} is whichever
task the interrupt happened to land on. Recording it as the waker produces a
false attribution and, worse, splices that bystander's unrelated history
into the chain. Detecting interrupt context requires the preemption count,
which on x86-64 is a per-CPU variable, and, because softirqs run on
interrupt exit after the hardirq bits are cleared, an explicit marker for
execution inside \texttt{irq\_exit\_rcu()} (\autoref{sec:chain}).

\subsection{eBPF}

eBPF lets user-supplied programs, verified for safety at load
time~\cite{gershuni2019verifier}, run at kernel hooks: tracepoints
(\texttt{tp\_btf}), function entry and exit (\texttt{fentry}/\texttt{fexit}),
and user-space probes (uprobes). Programs keep state in maps, including
\emph{task-local storage}, a per-task slot freed with the
task~\cite{bpf-task-storage}, and send records to user space through a ring
buffer~\cite{bpf-ringbuf}. CO-RE relocations let one object load on many
kernels~\cite{bpf-core}. The verifier imposes the constraints that shape
\sys's design: a 512-byte stack, bounded loops, a limit of one million
verified instructions, and constant-size memory copies.
